\documentclass[10pt,a4paper]{amsart}
\usepackage[margin=19mm]{geometry}
\usepackage{amsmath,amssymb,mathtools}
\usepackage[scr=rsfs,cal=euler]{mathalfa}
\usepackage{xfrac}
\usepackage[unicode,hidelinks,bookmarks=false]{hyperref}
\newcommand{\ie}{i.e.\ }
\newcommand{\cf}{cf.\ }
\newcommand{\resp}{respectively\ }
\newcommand{\Subsection}{\subsection}
\providecommand{\nobreakdash}{\nobreak\hbox{-}}
\setlength{\emergencystretch}{2em}
\multlinegap0pt
\usepackage{lmodern}
\usepackage{mathtools}
\usepackage{array}
\usepackage{xurl}
\usepackage{booktabs,tabularx}
\usepackage{enumitem}
\newcommand{\YZ}[1]{{#1}}
\newcommand{\GG}[1]{{#1}}
\newcommand{\VV}[1]{{#1}}
\hypersetup{
	colorlinks=true,
	linkcolor=blue,
	citecolor=blue,
	urlcolor=blue
}
\newtheorem{theorem}{Theorem}
\newtheorem{problem}[theorem]{Problem}
\newtheorem{proposition}[theorem]{Proposition}
\newtheorem{lemma}[theorem]{Lemma}
\newtheorem{corollary}[theorem]{Corollary}
\newtheorem{conjecture}[theorem]{Conjecture}
\newtheorem{claim}[theorem]{Claim}
\newtheorem{question}[theorem]{Question}
\newtheorem{obs}[theorem]{Observation}
\newtheorem{definition}[theorem]{Definition}
\newtheorem{remark}[theorem]{Remark}
\newcommand{\supp}{\operatorname{supp}}
\newcommand{\rank}{\operatorname{rank}}
\newcommand{\dotcupop}{\mathbin{\dot\cup}}
\newcommand{\foe}{f_{\mathrm{oe}}}
\newcommand{\hl}{h_{\ell}}
\newcommand{\fo}{f_o}
\newcommand{\V}{V}
\newcommand{\chiG}{\chi(G)}
\newcommand{\tw}{{\rm tw}}
\newcommand{\pcf}{\rm pfc}
\newcommand{\Ftwo}{\mathbb{F}_2}
\newcommand{\eps}{\varepsilon}
\newcommand{\E}{\mathbb E}
\DeclareMathOperator{\grad}{grad}
\newcommand{\C}{\mathcal C}
\newcommand{\B}{\mathcal B}
\newcommand{\F}{\mathbb{F}}
\newcommand{\im}{\operatorname{im}}
\newcommand{\1}{\mathbf{1}}
\DeclareMathOperator{\rk}{rank}
\DeclareMathOperator{\spn}{span}
\begin{document}
\title[Bounds on Odd and Odd-Even Induced Subgraphs]{Bounds on Odd and Odd-Even Induced Subgraphs}
\author{Qiwen Guo, Gregory Gutin, Yiming Hao, Yongtang Shi, Yong Zhang,; Yacong Zhou}
\date{}
\maketitle
\noindent\emph{Source and licence.} Bounds on Odd and Odd-Even Induced Subgraphs. Version arXiv:2608.02339v1 (CC BY 4.0). This material is distributed under the Creative Commons Attribution 4.0 International licence, http://creativecommons.org/licenses/by/4.0/, as recorded in the arXiv OAI licence metadata for this exact version and shown on the abstract page https://arxiv.org/abs/2608.02339v1. Full rights notice: licenses/arxiv-2608.02339-CC-BY-4.0.txt.\par\smallskip
\noindent\textit{Editorial scope.} Counted unit: the original \texttt{theorem} environment \texttt{thm:pointwise} at source lines 412--426 together with its complete proof at lines 427--468; the delivered document keeps source lines 209--469 so that every referenced label and every citation resolves inside it. Native editable source is the paper's own concatenated TeX; the AMS wrapper is editorial formatting only. No publisher class, upstream script or unrelated artwork is redistributed.
	\begin{lemma}\label{lem:completion}
		Let $B=(X,Y;F)$ be a bipartite graph and let $\ell:X\cup Y\to\F_2$ be a labelling such that
		$
		d_Y(x)\equiv \ell(x)\pmod 2
		\text{  for every }x\in X.
		$
		Then there are sets $X'\subseteq X$ and $Y'\subseteq Y$ with $|X'|+|Y'|\ge |X|$ such that
		$
		d_{Y'}(x)\equiv \ell(x)\pmod 2
		$
		for every $x\in X'$ and
		$
		d_{X'}(y)\equiv \ell(y)\pmod 2
		$
		for every $y\in Y'$.
	\end{lemma}
	
	\begin{proof}
		Let $A$ be the $X\times Y$ bipartite adjacency matrix of $B$, viewed over
		$\F_2$, and let $\ell_X$ denote the column vector of labels on $X$.  For $S\subseteq Y$ and
		$x\in X$ we have $(A\1_S)_x\equiv d_S(x)\pmod 2$. Thus, the assumption of this lemma translates to $A\mathbf 1_Y=\ell_X$.
		Let $Y'\subseteq Y$ be an inclusion-minimal set such that
		\begin{equation}\label{3.2}
			A\mathbf 1_{Y'}=\ell_X.
		\end{equation}
        Let $A_{Y'}$ be the submatrix formed by the columns indexed
		by $Y'$. If $\ell_X=0$, then $Y'=\emptyset$ and therefore we are done by setting $X':=X$ and $Y':=\emptyset$. 
		
		Thus, we may assume that $\ell_X\neq 0$ and therefore $Y'\neq \emptyset$. Then, the columns of $A_{Y'}$ are linearly independent, i.e., $\rank(A_{Y'})=|Y'|$. Indeed, if the columns indexed by a nonempty set $Y''\subseteq Y'$ summed to zero (therefore $Y''\subsetneq Y'$ as $\ell_X\neq 0$), then
		\[
		A\mathbf 1_{Y'\setminus Y''}
		=A\mathbf 1_{Y'}-A\mathbf 1_{Y''}
		=\ell_X,
		\]
		contrary to the minimality of $Y'$.  
		
		Let $\ell_{Y'}$ be the vector of labels on $Y'$.  Since $A_{Y'}^{\top}$ has full row
		rank $|Y'|$, the system
		\begin{equation}\label{3.4}
			A_{Y'}^{\top} z=\ell_{Y'} 
		\end{equation}
		has a solution $z_0\in\F_2^X$.  Its solution set is therefore $z_0+\ker A_{Y'}^{\top}$. In addition, by rank-nullity, we have
	\[
			\dim \ker A_{Y'}^{\top}=|X|-\rank (\ker A_{Y'}^{\top})=|X|-|Y'| . \]
		Let $d=|X|-|Y'|$.  If $d=0$, set $I:=\emptyset$ and $z:=z_0$.  Otherwise let $\psi_1,\dots,\psi_d$ be a basis of $\ker A_{Y'}^{\top}$ and let $A'$ be the $d\times |X|$ matrix with rows $\psi_1,\dots,\psi_d$.  As row rank equals column rank, $A'$ has $d$ linearly independent columns. Let $I\subseteq X$ be the set of the corresponding coordinates, so that the $d\times d$ submatrix $A''$ of $A'$ formed by these columns is invertible.  Hence the restriction map $\ker A_{Y'}^{\top}\to\F_2^{I}$, $\psi\mapsto \psi|_I$, is a bijection, and in particular there is $w\in \ker A_{Y'}^{\top}$ with $w|_I=\1_I+z_0|_I$.  Let $z:=z_0+w$. Then, $z$ solves~\eqref{3.4} and $z|_I=\1_I$.
		
		Let $X'=\supp (z)$.  In either case $I\subseteq X'$, and therefore
		\[
			|X'|\ge |I|=|X|-|Y'|, 
		\]
		so that $|X'|+|Y'|\ge |X|$.  Finally, for $x\in X'$, equation~\eqref{3.2} gives
		$
		d_{Y'}(x)\equiv \ell(x)\pmod 2,
		$
		and for $y\in Y'$, the coordinate of~\eqref{3.4} indexed by $y$ gives
		$
		d_{X'}(y)\equiv \ell(y)\pmod 2.
		$
	\end{proof}
	
	\section{Labelled Cuts}\label{sec:labelled}
	
	We first generalize Zeng's odd-cut notion, which is a special case of our notion below when $\ell\equiv 1$. Let $G$ be a graph with a labelling $\ell:V(G)\to\F_2$.  A set $S\subseteq V(G)$ is an \emph{$\ell$-cut set} if it has a partition $S=P\dotcupop Q$ such that $d_Q(v) \equiv \ell(v)\pmod 2 \mbox{ for all } v\in P \mbox{ and } d_P(v) \equiv \ell(v)\pmod 2 \mbox{ for all } v\in Q$. The two parts are allowed to be empty.
	
	
	We use the following classical theorem (see, for example,
	Lov\'asz~\cite[Problem~5.17]{Lovasz}).
	
	\begin{theorem}[Gallai's even-partition theorem]\label{thm:gallai}
		Every graph $H$ has a partition $V(H)=U\dotcupop W$ such that both $H[U]$
		and $H[W]$ have all degrees even.
	\end{theorem}
	
	For $\ell\equiv 1$, Theorem~\ref{thm:labelled-cut} below is Zeng's odd-cut theorem
	\cite[Theorem~1.2]{Zeng}, and its proof is a modification of the proof of \cite[Theorem~1.2]{Zeng}.
	We provide the proof for the sake of paper's self-containment. 
	
	\begin{theorem}\label{thm:labelled-cut}
		Let $G$ be a graph with a labelling $\ell:V(G)\to\F_2$, and let $S$ be an
		$\ell$-cut set.  Then there is a partition
		$
		S=S_0\dotcupop S_1
		$
		such that both $S_0$ and $S_1$ are $\ell$-admissible. 
	\end{theorem}
	
	\begin{proof}
		Choose a witnessing partition $S=P\dotcupop Q$.  Apply
		Theorem~\ref{thm:gallai} to $G[S]$ and obtain a partition
		$S=U\dotcupop W$ such that $G[U]$ and $G[W]$ are even.  Set
		$S_0=(U\cap P)\cup(W\cap Q)$
		and 
		$ S_1=(W\cap P)\cup(U\cap Q).$
		Clearly, $S=S_0\dotcupop S_1$.
		
		Fix $i\in\{0,1\}$ and $v\in S_i$.  Let $A\in\{U,W\}$ be the part containing
		$v$, and let $B\in\{P,Q\}$ be the part not containing $v$.  By the
		definition of $S_i$,
		$
		S_i=A\mathbin{\triangle}B,
		$
		where $\triangle$ denotes symmetric difference. 
		%Since \(S=P\mathbin{\dot\cup}Q=U\mathbin{\dot\cup}W\), we have
		%\[
		%\begin{aligned}
		%U\triangle Q=W\triangle P
		%  &=(U\cap P)\cup(W\cap Q)=S_0,\\
		%W\triangle Q=U\triangle P
		%   &=(W\cap P)\cup(U\cap Q)=S_1.
		%\end{aligned}
		%\]
		%Indeed, these identities follow directly from
		%\(C\triangle D=(C\cap(S\setminus D))\cup((S\setminus C)\cap D)\).
		%Now, if \(v\in S_0\), then either \(v\in U\cap P\), in which case
		%\((A,B)=(U,Q)\), or \(v\in W\cap Q\), in which case
		%\((A,B)=(W,P)\). Thus \(A\triangle B=S_0\).
		%Similarly, if \(v\in S_1\), then either
		%\((A,B)=(W,Q)\) or \((A,B)=(U,P)\), and consequently
		%\(A\triangle B=S_1\). Hence \(S_i=A\triangle B\).
		Therefore, modulo $2$,
		$
		d_{S_i}(v,G)
		\equiv d_A(v,G)+d_B(v,G)
		\equiv 0+\ell(v)
		\equiv \ell(v).
		$
		The first term is even because $G[A]$ is even, and the second has parity
		$\ell(v)$ because $P\dotcupop Q$ witnesses that $S$ is an $\ell$-cut set.
		Thus both $S_0$ and $S_1$ are $\ell$-admissible.  \end{proof}
	
	
	
	The labelled odd-cut theorem is useful only after a large $\ell$-cut set has been found. For arbitrary labels, imposing the correct cross parity on both parts at once is inconvenient, and this is where Lemma~\ref{lem:completion} is applied: given disjoint $P,Q\subseteq V(G)$ such that every $q\in Q$ has the prescribed parity of degree into $P$, apply the lemma to the
    bipartite graph $G[P,Q]$, whose edges are the edges of $G$ between $P$ and $Q$, with $X:=Q$ and $Y:=P$, to obtain the resulting sets $X'\subseteq Q$ and $Y'\subseteq P$ \VV{that} satisfy $|X'|+|Y'|\ge |Q|$ and are the two parts of an $\ell$-cut set $Y'\dotcupop X'$ of $G$.  Theorem~\ref{thm:labelled-cut} then gives $h_\ell(G)\ge (|X'|+|Y'|)/2\ge |Q|/2$, which is the following corollary.
	
	\begin{corollary}\label{cor:one-sided}
		Let $P,Q\subseteq V(G)$ be disjoint and suppose that
		$
		d_P(q,G)\equiv \ell(q)\pmod 2
		\mbox{ for all } q\in Q.
		$
		Then
		$
		h_\ell(G)\ge |Q|/2.
		$
	\end{corollary}
	%\begin{proof}
	%Apply Lemma~\ref{lem:completion} to the bipartite graph $G[P,Q]$.  It gives
	%$C\subseteq P$ and $R\subseteq Q$ such that $C\cup R$, with shores $C$ and
	%$R$, is an $\ell$-cut set of $G$ and $ |C|+|R|\ge |Q|.$
	%Theorem~\ref{thm:labelled-cut} now gives
	%$$
	% h_\ell(G)\ge \frac{|C|+|R|}{2}\ge \frac{|Q|}{2}.
	%$$
	%\end{proof}
	
	%\begin{remark}[Why two stages are needed]
	%Lemma~\ref{lem:completion} controls degrees only in the bipartite graph
	%$G[C,R]$; edges inside $C$ or inside $R$ may change the degrees in the
	%induced subgraph $G[C\cup R]$.  Thus the completion lemma alone does not
	%finish the problem in the original graph.  Its output is instead a labelled
	%cut, and Theorem~\ref{thm:labelled-cut} absorbs all same-shore edges by the
	%Gallai switching argument.  This two-stage synthesis---linear completion
	%followed by labelled odd-cut switching---is the essential structural
	%modification.
	%\end{remark}
	
	Let $\gamma(G)$ denote the domination number of $G$.  We use the following
	result of Bollob\'as and Cockayne~\cite{BC}; see also~\cite{Haynes}.
	
	\begin{theorem}[Bollob\'as--Cockayne]\label{thm:BC}
		Every graph $G$ without isolated vertices has a minimum dominating set
		$D$ such that, for each $u\in D$, there exists a vertex
		$q_u\in V(G)\setminus D$ satisfying
		\[
		N_G(q_u)\cap D=\{u\}.
		\]
	\end{theorem}
	
	The vertices $q_u$ are external private neighbours and are necessarily
	distinct.
	
	\begin{proposition}\label{prop:gamma}
		Let $G$ have no isolated vertices and let $\ell:V(G)\to\F_2$.  Then
		$ h_\ell(G)\ge \gamma(G)/2.$
	\end{proposition}
	\begin{proof}
		Choose a minimum dominating set $D$ as in Theorem~\ref{thm:BC}, and choose
		one external private neighbour $q_u$ for each $u\in D$.  Put
		$
		Q:=\{q_u:u\in D\}
		$
		and
		$
		P:=\{u\in D:\ell(q_u)=1\}.
		$
		For every $u\in D$, the private-neighbour property gives
		$d_P(q_u,G)=\ell(q_u).$ Thus, $d_P(q_u,G)\equiv\ell(q_u)\pmod2$ for all $q_u\in Q$.  Since
		$|Q|=|D|=\gamma(G)$, Corollary~\ref{cor:one-sided} yields the result.
	\end{proof}
	
		Note that in the all-odd case, Zeng takes all of $D$ and one private neighbour for each $u\in D$. The resulting $2\gamma(G)$ vertices form an odd-cut set and give an odd induced subgraph of order at least $\gamma(G)$.  Under an arbitrary prescription, taking all of $D$ no longer gives the correct cross parity at a private neighbour labelled $0$. The price of arbitrary labels is that the guaranteed contribution is $\gamma(G)/2$ rather than $\gamma(G)$. At the end, this leads to $n/6$ rather than $n/5.$


	\begin{theorem}\label{thm:pointwise}
		Let $G$ be an $n$-vertex graph without isolated vertices and let
		$\ell:V(G)\to\F_2$.  Then
		\[
		h_\ell(G)\ge
		\max\left\{
		\frac{\gamma(G)}2,
		\frac{n-\gamma(G)}4
		\right\}.
		\]
		In particular,
		\VV{$
		h_\ell(G)\ge n/6,
		$ and therefore $\foe(G)\ge n/6$.}
	\end{theorem}

	\begin{proof}
     By \VV{Proposition}~\ref{prop:gamma}, we only need to show $
	h_\ell(G)\ge \frac{n-\gamma(G)}{4}$. 
		
		\begin{claim}\label{cl:n-gamma/4}
			$h_\ell(G)\ge \frac{n-\gamma(G)}{4}$. 
		\end{claim}
		
		\begin{proof}
					Let $D$ be any minimum dominating set, so $|D|=\gamma(G)$.  Choose a random
			set $Z\subseteq D$ by including every vertex independently with probability
			$1/2$, and $Q_Z:=\bigl\{v\in V(G)\setminus D: d_Z(v,G)\equiv\ell(v)\pmod2\bigr\}.$
			
			Since $D$ dominates $G$, every $v\in V(G)\setminus D$ has a neighbour in $D$. \VV{Fix such a neighbour $x_v\in D$; after all the choices for the vertices of $D\setminus\{x_v\}$ have been exposed, exactly one of the two possible choices for $x_v$ makes $d_Z(v,G)$ congruent to $\ell(v)$ modulo~2.} Thus, we have for every $v\in V(G)\setminus D$, $\Pr(v\in Q_Z)=\frac12$.
			By linearity of expectation,
			$
			\mathbb E|Q_Z|=\frac{n-\gamma(G)}2,
			$
			and therefore some choice of $Z$ satisfies
			\[
			|Q_Z|\ge \frac{n-\gamma(G)}2.
			\]
			For that choice, the pair $(Z,Q_Z)$ satisfies the hypothesis of
			Corollary~\ref{cor:one-sided} (with $P:=Z$).  Therefore
			\[
			h_\ell(G)
			\ge \frac{|Q_Z|}{2}
			\ge \frac{n-\gamma(G)}4,
			\]
		completing the proof of the claim.
		\end{proof}
		
		Thus, combining Proposition \ref{prop:gamma} and Claim \ref{cl:n-gamma/4}, we have
		\[
		h_\ell(G)\ge
		\max\left\{
		\frac{\gamma(G)}2,
		\frac{n-\gamma(G)}4
		\right\}\geq \frac{1}{3}\cdot \frac{\gamma(G)}{2}+\frac{2}{3}\cdot\frac{n-\gamma(G)}{4}=\frac{n}{6},
		\]
		completing the proof.
	\end{proof}
	

\begin{thebibliography}{99}
\bibitem[BC]{BC} B.~Bollob\'as and E.~J. Cockayne, Graph-theoretic parameters concerning domination, independence, and irredundance, J. Graph Theory 3 (1979), 241--249.
\bibitem[Haynes]{Haynes} T.~W. Haynes, S.~T. Hedetniemi and M.~A. Henning, \emph{Domination in Graphs: Core Concepts}, Springer, Cham, 2023.
\bibitem[Lovasz]{Lovasz} L.~Lov\'asz, \emph{Combinatorial Problems and Exercises}, 2nd ed., AMS Chelsea Publishing, Providence, RI, 1993.
\bibitem[Zeng]{Zeng} Q.~Zeng, Large odd induced subgraphs via odd cuts, arXiv: 2607.23266v1, 2026.
\end{thebibliography}
\end{document}
